% =====================================================================
\part{O que sustenta peso de conclusão, e o que sustenta peso de indício}
% =====================================================================

\chapter{A escala de graus de estabelecimento}

O documento traz, na Discussão, um quadro que atribui a cada frente um grau e
diz o que elevaria esse grau. Esse quadro é a sua melhor defesa contra a
pergunta ``e o que o senhor realmente provou?'', porque a resposta já está
escrita e graduada.

\vspace{6pt}
\begin{longtable}{@{}p{2.0cm}p{5.4cm}p{2.4cm}p{4.6cm}@{}}
\toprule
\textbf{Frente} & \textbf{Afirmação} & \textbf{Grau} & \textbf{O que elevaria} \\
\midrule
\endfirsthead
\toprule
\textbf{Frente} & \textbf{Afirmação} & \textbf{Grau} & \textbf{O que elevaria} \\
\midrule
\endhead
1. O padrão &
Conversão e rebanho deslocam-se ao norte; o gradiente latitudinal persiste &
\textbf{estabelecido} &
já sobrevive a pixel, malha, classe e barra de erro; o que resta é a acurácia do
próprio classificador \\
\addlinespace
2. Mecanismo &
Duas lógicas coexistem em todo o estado; a região explica pouco da idade da
pastagem &
\textbf{estabelecido} &
o que a região não explica pede dado por imóvel rural \\
\addlinespace
3a. Deslocamento &
O canal local e contemporâneo não apresenta a assinatura exigida &
\textbf{nulo}, com poder declarado no eixo temporal &
desenho de maior alcance espacial e defasagem longa; travessia da divisa
estadual \\
\addlinespace
3b. \emph{Drive} comum &
A resposta ao choque distribui-se ao longo do eixo Sul--Norte, no sentido
previsto e aquém do corte; nem a força nem a característica do eixo são
identificadas &
\textbf{corroborante} &
\emph{shifter} com variação transversal, janela maior ou painel multiestadual;
exposições que não se ordenem com a latitude \\
\addlinespace
4. Teto de oferta &
A freada do Sul é compatível com atrito de oferta, e não com esfriamento da
demanda &
\textbf{eliminação \emph{e} mecanismo medido} &
cadastro ambiental pixel a pixel; e o Ato~III é curto \\
\bottomrule
\end{longtable}

\vspace{12pt}
\begin{nabanca}[Como usar este quadro numa resposta]
``Eu classifico as minhas quatro frentes em graus diferentes, e o quadro da
Discussão diz qual é qual e o que elevaria cada um. As duas primeiras eu
considero estabelecidas. A terceira tem duas metades: um nulo com poder
declarado, e uma leitura corroborante que eu não apresento como demonstrada. A
quarta se sustenta por eliminação e por um mecanismo medido. Se a pergunta é
onde eu me arrisco menos, é nas duas primeiras; onde eu me arrisco mais, é na
3b, e é por isso que ela vem com o nome que vem.''
\end{nabanca}

\section{Os limites que atravessam todas as frentes}

Saber enunciá-los espontaneamente é o que transforma uma arguição em conversa.

\vspace{4pt}
\textbf{O alcance é intraestadual.} Tudo o que foi testado ocorre dentro de
Goiás. Soja goiana deslocando pecuária para o Pará ou para o Maranhão é outra
pergunta, e os dados aqui reunidos não a fazem.

\textbf{A medida depende de um classificador, e ele mudou de comportamento no
fim da série.} O intervalo entre as duas réguas \emph{contém} o artefato; ele
não o corrige.

\textbf{O terceiro ato tem cinco anos.} A freada do Sul, a mudança de endereço da
fronteira e a estabilização do estoque meridional são sinal inicial, e não regime
consolidado. A estabilização pode ainda revelar-se uma pausa, e o desenho não
distingue as duas hipóteses; só o tempo o fará.

\textbf{``Convertível'' e ``protegida'' são aproximações declaradas.} São tetos
adequados ao argumento de estoque, mas medem exposição por aproximação, e não
lote a lote.

\textbf{O recorte por mesorregião é grosso.} Cinco faixas para um estado do
tamanho de Goiás escondem variação interna, e onde foi possível descer à
malha fina (as 166 AMCs, e até o pixel), a conclusão não mudou.

\textbf{Sobre a idade da pastagem, nada se afirma quanto a tendência.} O eixo do
tempo está comprometido nas duas pontas: antes, pela idade truncada no início da
série; depois, pela mudança de rótulo.

\chapter{As sete autocorreções, e como transformá-las em ativo}

Este é, provavelmente, o capítulo mais importante do caderno para a impressão
que você deixa. Um trabalho que registra sete resultados que teriam sido
publicados e foram superados por um teste adicional está dizendo à banca algo
que nenhuma quantidade de robustez adicional diria.

\begin{quote}\itshape\color{cortinta}
Nenhum deles era ingênuo: em cada caso, a versão superada era a leitura mais
plausível com o que havia até ali. Em quatro casos ela afirmava \emph{mais} do
que a versão que ficou no lugar, em dois contrariava a tese, e no último a
versão corrigida sustenta \emph{mais} do que a superada.
\end{quote}

\section{As sete, com a resposta de uma frase para cada}

\vspace{4pt}
\textbf{1. ``A agricultura desacelerou depois de 2019.''} \onde{D25--D26} \\
O que o teste mostrou: a classe parou de registrar, e as medidas de campo
marcham na mesma janela. Consequência: toda medida de agricultura que alcance
2020--2024 passa a ser um intervalo.

\textbf{2. ``O Norte antecede o Sul no tempo.''} \onde{D16} \\
O que o teste mostrou: regressão espúria. Sob método robusto à integração, a
precedência some \emph{nas duas direções}. Consequência: o veredito é
\emph{sem líder}, e nenhum resultado de precedência em série agregada é lido
como causal sem diagnóstico de estacionariedade, método robusto e bateria de
placebos.

\textbf{3. ``O \emph{drive} comum é estatisticamente significante.''} \onde{D20} \\
O que o teste mostrou: a inferência apropriada ao desenho leva o $p$ de cerca de
0,03 para 0,07--0,13. Consequência: o veredito vira \emph{corroborante, não
estabelecido}, e os p-valores agrupados de 0,026 e 0,031 \textbf{não devem
ser citados como significância}.

\textbf{4. ``A marcha ao norte'', sem barra de erro.} \onde{D19} \\
O que o teste mostrou: com IC por \emph{bootstrap}, a vegetação natural passa a
``ancorada''. Consequência: um deslocamento cujo intervalo inclui zero nunca é
reportado em quilômetros.

\textbf{5. ``O gradiente do \emph{drive} comum é o de aptidão
edafoclimática.''} \onde{D28} \\
O que o teste mostrou: posta a latitude na mesma regressão, a aptidão perde 62\%
da magnitude e a significância nas duas réguas. Consequência: o resultado passa
a ser enunciado como gradiente \textbf{Sul--Norte}, e a aptidão é descrita como
a régua com que ele foi medido, não como o canal identificado.

\textbf{6. ``A taxa de conversão não cai com o esgotamento do estoque.''}
\onde{D29} \\
O que o teste mostrou: o nulo vinha de 14\% do painel com esgotamento fora do
domínio (até $-85$ numa fração que deveria viver entre 0 e 1). Tratado o
defeito, a taxa \textbf{cai}. Consequência: a Perna 4 sai \emph{fortalecida} ---
é o único caso em que a correção sustenta mais do que a versão superada.

\textbf{7. ``A formação florestal emite mais que a savânica, apesar de perder
menos área.''} \onde{D30} \\
O que o teste mostrou: a afirmação dependia de uma \emph{razão} entre duas
densidades superar 2,64, e os três cenários de sensibilidade moviam as duas
juntas. Com os estoques do inventário oficial a razão é 1,57, e a savânica
responde por 573 Mt contra 340. Consequência: quando uma conclusão depende de
uma razão entre parâmetros, é a razão que precisa entrar na sensibilidade.

\vspace{8pt}
A essas somam-se três apresentadas dentro das frentes correspondentes: o
gradiente latitudinal de idade da pastagem, o plantio direto como mecanismo, e a
leitura da vegetação como bloco homogêneo ao norte.

\section{Como narrá-las, e como não narrá-las}

Há um jeito errado e um jeito certo, e a diferença é de foco.

\begin{armadilha}[O jeito errado]
Narrar como ``crença desfeita'': ``eu \emph{achava} que a aptidão era o canal, e
descobri que não''. Isso convida a banca a perguntar o que mais você acha e
ainda não testou.
\end{armadilha}

\begin{nabanca}[O jeito certo]
Narrar como \textbf{resultado superado} --- isto é, como o que \emph{teria sido
publicado} se o teste adicional não tivesse sido feito: ``a versão que teria ido
para o papel dizia X; o teste Y mostrou que ela dependia de Z; o que ficou no
lugar é W''. O foco sai da sua crença e vai para o teste, que é onde está o
mérito.
\end{nabanca}

\begin{nabanca}[E a moldura geral, para dizer uma vez só]
``Eu mantenho um registro de resultados superados porque um trabalho que só
mostra o que sobreviveu não permite avaliar o filtro. Sete deles estão no
quadro da §4.6, cada um com o teste específico que o derrubou e a decisão
numerada que ficou no lugar, para impedir que a escolha mais simples seja
restabelecida por conveniência --- inclusive por mim, no futuro.''
\end{nabanca}

\chapter{Frases proibidas e as suas substitutas}
\label{cap:frases}

Cada linha desta tabela é um excesso de afirmação que soaria bem e não se
sustenta. A coluna da direita é o que dizer no lugar.

\vspace{8pt}
% longtable, e não tabular: com quinze linhas altas ela não cabe numa página, e
% um tabular que não cabe é empurrado inteiro para a página seguinte — deixando
% atrás dele uma página em branco sob o título do capítulo.
\begin{longtable}{@{}p{6.5cm}p{8.2cm}@{}}
\toprule
\textbf{Não dizer} & \textbf{Dizer} \\
\midrule
\endfirsthead
\toprule
\textbf{Não dizer} & \textbf{Dizer} \\
\midrule
\endhead
``O iLUC foi refutado'' / ``o deslocamento indireto não existe'' &
``A assinatura testada não aparece nos dados examinados, com o poder disponível.
O canal de longo alcance e o de defasagem longa ficam fora do teste'' \\
\addlinespace
``O câmbio causa a expansão da fronteira'' &
``O nível causal do canal cambial vem da literatura, em escala nacional. O que
eu estimo é a \emph{geografia} desse efeito dentro do estado'' \\
\addlinespace
``O efeito é significativo ($p = 0{,}026$)'' &
``Sob a inferência apropriada ao desenho, o $p$ fica entre 0,07 e 0,13. O que
sustenta a leitura é a especificidade, não a significância'' \\
\addlinespace
``A aptidão do solo explica o gradiente'' &
``A aptidão é a régua com que o eixo Sul--Norte foi medido. Qual característica
do eixo faz a diferença, este desenho não separa'' \\
\addlinespace
``Restam 6,56 Mha \emph{disponíveis} para expansão'' &
``Restam 6,56 Mha de vegetação natural ainda \emph{exposta} à conversão. É o
Cerrado que ainda pode ser perdido'' \\
\addlinespace
``O Cerrado do Sul está se regenerando'' &
``O fluxo reverso é quase todo dirigido à formação savânica e é aproximadamente
simétrico: é oscilação de classificador, não recuperação'' \\
\addlinespace
``A vegetação natural andou 7,6 km ao norte'' &
``A vegetação natural está \emph{ancorada}: o intervalo inclui zero'' \\
\addlinespace
``O Sul preservou'' &
``Restou pouco Cerrado em pé a ser suprimido, e o que restou é majoritariamente
a fisionomia que a fronteira não consome'' \\
\addlinespace
``O que resta está desprotegido'' &
``Está fora do grupo de Proteção Integral, que é a única barreira de que tenho
medida espacial completa. Reserva Legal e APP seguem incidindo, e eu não as
meço'' \\
\addlinespace
``A união de Agricultura com Mosaico é a medida correta'' &
``A união é o teto do intervalo: ela superconta tanto quanto a classe isolada
subconta, e responde a uma pergunta mais grossa'' \\
\addlinespace
``O Código Florestal não teve efeito'' &
``Ele não produz quebra empírica em nenhuma das três séries, o que rebaixa o seu
estatuto de marco periodizador, e não diz nada sobre efeitos por outras
vias'' \\
\addlinespace
``O modelo de diferenças em diferenças mostra o efeito da política'' &
``Rebaixei esse desenho: os marcos são federais e não deixam grupo não tratado.
O que o coeficiente mede é exposição diferencial ao mesmo choque'' \\
\addlinespace
``A fronteira parou'' &
``A fronteira mudou de endereço. O fluxo bruto estadual é praticamente o mesmo
entre os Atos II e III'' \\
\addlinespace
``A abertura de área trouxe desenvolvimento'' (ou o inverso, como efeito) &
``No período e na escala observados, a expansão da área não aparece
\emph{associada} a convergência de desenvolvimento. A leitura é associativa'' \\
\addlinespace
``A taxa cai porque o que restou é mais difícil'' &
``Caem juntos o estoque e a taxa. É regularidade compatível com atrito de
oferta, e não demonstração de que a dificuldade seja a causa'' \\
\bottomrule
\end{longtable}

\vspace{12pt}
\begin{nabanca}[A regra que gera todas as linhas acima]
Toda afirmação forte do trabalho vem acompanhada da lista dos seus limites, no
mesmo lugar. Se você se pegar dizendo uma frase forte sem a ressalva ao lado, a
frase provavelmente não é do trabalho: é a versão dele que soa melhor.
\end{nabanca}
